\chapter*{Как читать эту книгу}
\addcontentsline{toc}{chapter}{Как читать эту книгу}

Эта книга объясняет одну модель: как при остывании трубы из циркониевого сплава в ней
выпадают пластинки гидрида и почему они ложатся вдоль окружности или поперёк неё, по радиусу.
Модель собиралась по шагам в течение нескольких сессий работы; здесь она изложена целиком,
от физики до кода, в том виде, в каком она есть сейчас, вместе с тем, что в ней не работает.

\paragraph{Три слоя.} Каждая глава построена одинаково.
\begin{enumerate}
\item \textbf{Картинка и интуиция.} Сначала --- что происходит, на пальцах и на рисунке.
  Синие блоки «Интуиция» можно читать отдельно от остального текста: из них складывается
  короткий рассказ о модели.
\item \textbf{Уравнения.} Затем --- формулы, откуда они берутся и что в них означает каждая буква.
  Главные формулы обведены золотой рамкой.
\item \textbf{Упражнения.} Зелёные блоки --- задачи, которые стоит решить на бумаге. Большинство
  решается за 5--15 минут с калькулятором; ответы и решения --- в приложении~\ref{app:answers}.
  Это не проверка, а способ «потрогать» числа: после них станет понятно, почему градус
  температуры важнее сотни мегапаскалей и почему пластичность нельзя выбросить.
\end{enumerate}
Красные блоки «Ограничение» честно говорят, где модель упрощает или где она пока не сходится
с опытом. Серые блоки «Где в коде» указывают файлы в папке \code{hydrides/automaton}.

\paragraph{Что нужно знать заранее.} Производные и интегралы, экспонента и логарифм, немного
линейной алгебры (матрицы $2\times2$, собственные векторы). Тензоры напряжений и деформаций
вводятся в главе~\ref{ch:misfit} с нуля, в объёме, нужном для модели. Преобразование Фурье
объясняется там же на одном примере.

\paragraph{Карта.} Рисунок~\ref{fig:F02_Map} --- оглавление в картинке: какие блоки есть в модели,
что они передают друг другу и в какой главе разобран каждый. Модель принимает на вход
\emph{условия опыта} (сплав, текстуру, водород, температурную историю, нагрузку) и выдаёт
\emph{структуру} гидридов и её меры, которые сравниваются с шлифами.

\fig{F02_Map}{Модель целиком. Внутри клеточного автомата четыре физических блока обмениваются
данными на каждом шаге по времени; новые пластинки меняют и поле напряжений, и водород вокруг.}

\paragraph{Обозначения.} Окружное направление трубы обозначается $\theta$ (оно же TD в
металлографии), радиальное --- $r$ (ND), осевое --- $z$ (RD). На всех рисунках
\textcolor{circ}{окружные} гидриды синие, \textcolor{rad}{радиальные} --- красные, нагрузка
\textcolor{load}{золотая}, пластичность \textcolor{plast}{фиолетовая}, водород \textcolor{hydr}{зелёный}.
Полный список обозначений и параметров --- в приложении~\ref{app:ref}.
